\chapter*{Introduction générale}
\addcontentsline{toc}{chapter}{Introduction générale}
\markboth{Introduction générale}{Introduction générale}

Le secteur du \textbf{tourisme} s'appuie de plus en plus sur le \textbf{numérique} pour informer, orienter et fidéliser les visiteurs. Les voyageurs consultent des avis, des cartes et des applications mobiles ; pourtant, l'information disponible est souvent \textbf{hétérogène}, \textbf{peu actualisée} ou \textbf{difficile à vérifier} sur le terrain.

Le projet \textbf{MoroccoCheck} répond à ce constat en proposant une \textbf{plateforme touristique communautaire} centrée sur le Maroc. L'application combine la consultation de \textbf{sites touristiques}, la \textbf{vérification de présence} par \textbf{géolocalisation} (check-in GPS avec contrôle de distance), la publication d'\textbf{avis} et une \textbf{gamification} (points, badges, classement). Un rôle \textbf{professionnel} permet aux gestionnaires de lieux de revendiquer des fiches et d'interagir avec les retours visiteurs, tandis que les \textbf{administrateurs} assurent la \textbf{modération} et le pilotage via une interface web dédiée.

\section{Problématique}

La question centrale de ce mémoire est la suivante :

\begin{quote}
\textit{Comment concevoir une plateforme touristique communautaire capable d'améliorer la fiabilité des informations grâce à la géolocalisation, à la participation des utilisateurs et à la modération, tout en restant évolutive sur le plan technique ?}
\end{quote}

\section{Objectifs du projet}

Les objectifs principaux sont :
\begin{itemize}[leftmargin=*]
  \item fournir un \textbf{catalogue consultable} de sites (liste, détail, catégories, carte) ;
  \item mettre en œuvre des \textbf{check-ins validés par GPS} (distance, précision, règles métier) ;
  \item permettre la \textbf{contribution communautaire} (avis, notes, médias selon les parcours) ;
  \item encourager l'engagement par la \textbf{gamification} ;
  \item gérer les \textbf{rôles} (touriste, contributeur, professionnel, administrateur) et les demandes d'évolution de statut ;
  \item offrir un \textbf{back-office} pour statistiques et modération.
\end{itemize}

\section{Méthodologie et structure du document}

Le travail s'appuie sur une analyse des besoins, une phase de conception (modélisation UML, base de données), puis une réalisation incrémentale (API, mobile, administration) avec tests et intégration continue.

Ce rapport est organisé en \textbf{trois chapitres} : le \textbf{cadre général et l'analyse des besoins}, la \textbf{conception de la solution}, puis la \textbf{réalisation et l'implémentation}. Une \textbf{conclusion générale}, un chapitre sur les \textbf{outils de développement} et une \textbf{bibliographie} complètent le document.
